\subsection{Adaptive Test-Time Sampling for Transolver}
\label{sec:adaptive_test_time_sampling}

We next construct the inference-time sampling measure from geometry, a draft prediction, or intermediate backbone features.

Let $r_{\psi}:\Omega\to[0,\infty)$ be a nonnegative refinement score, where $\psi$ denotes the information used to construct it, and define
\begin{equation}
    \nu_{\psi}(x)
    :=\frac{r_{\psi}(x)}{\int_{\Omega}r_{\psi}(y)\,\mathrm{d}\rho(y)},
    \qquad
    \widetilde{\rho}_{\gamma,\psi}
    :=(1-\gamma)\rho+\gamma(\nu_{\psi}*\rho),
    \qquad \gamma\in[0,1].
    \label{eq:unified_refinement_measure}
\end{equation}
Here $\nu_{\psi}*\rho$ denotes the reweighted empirical measure induced by the refinement score. The correction used in the point-to-slice aggregation is
\begin{equation}
    c_{\gamma,\psi}(x)
    :=\frac{\mathrm{d}\rho}{\mathrm{d}\widetilde{\rho}_{\gamma,\psi}}(x)
    =\frac{1}{(1-\gamma)+\gamma\nu_{\psi}(x)}.
    \label{eq:unified_refinement_correction}
\end{equation}

\subsubsection{Physics-Informed Density Refinement}
\label{sec:physics_informed_refinement}

The first policy constructs $r_{\psi}$ from prior information about the geometry
or the physical problem. Let $s_{\mathrm{prior}}(x)\geq0$ be a problem-dependent
indicator whose large values identify regions expected to require finer resolution.
This indicator can be obtained from a known geometric feature, a numerical
refinement rule, or a physical quantity available before inference. We set
\begin{equation}
    \nu_{\mathrm{prior}}(x)
    :=\frac{s_{\mathrm{prior}}(x)}
    {\int_{\Omega}s_{\mathrm{prior}}(y)\,\mathrm{d}\rho(y)}
    \label{eq:prior_refinement_kernel}
\end{equation}
and use it in \eqref{eq:unified_refinement_measure} to draw the inference samples. This policy does not require a preliminary network evaluation; once the prior indicator is available, points are sampled from $\widetilde{\rho}_{\gamma,\mathrm{prior}}$ and inferred with the correction in \eqref{eq:unified_refinement_correction}.

\subsubsection{Coarse-Prediction-Guided Density Refinement}
\label{sec:coarse_prediction_refinement}

This strategy is motivated by PDE-Refiner~\citep{lippe2023pderefiner}
and ATCA~\citep{jenkins2026atca}. When no reliable prior is available,
we draw a small coarse sample $S_0$ from $\rho$ and evaluate the frozen
operator:
\begin{equation}
    \widehat{u}^{(0)}
    :=\mathcal{G}_{\theta}\!\left(a\big|_{S_0}\right),
    \qquad |S_0|\ll N,
    \label{eq:coarse_draft_prediction}
\end{equation}
where $a\big|_{S_0}$ denotes the input field restricted to $S_0$.
We then construct a nonnegative refinement score from the coarse prediction,
for example,
\begin{equation}
    s_{\mathrm{coarse}}(x)
    :=\Phi\!\left(
        \widehat{u}^{(0)}(x),
        \nabla\widehat{u}^{(0)}(x),
        \Delta\widehat{u}^{(0)}(x),
        \ldots
    \right),
    \qquad \Phi\geq0,
    \label{eq:coarse_refinement_score}
\end{equation}
where $\Phi$ denotes a task-dependent combination of coarse predictive
features. These quantities are estimated from the local discrete neighborhood
and extended to candidate locations when necessary. The resulting score is
normalized through \eqref{eq:unified_refinement_measure} to define
$\widetilde{\rho}_{\gamma,\mathrm{coarse}}$.

The resulting density concentrates samples in regions that the coarse prediction
identifies as spatially informative or task-relevant.

\subsubsection{Cache-Aware Learned Refinement}
\label{sec:cache_aware_refinement}

The coarse prediction provides features correlated with local physical or geometric variation, but the subsequent computation of the first pass is otherwise discarded. We reuse it as follows. Let $S_1=S_0\cup B$ be the final point set, where $B$ contains the newly allocated points. For the first Physics-Attention layer, let $h_i^{\mathrm{in}}$ be its pointwise input feature and $w_{i,j}^{(1)}$ its slice weight. Define
\begin{equation}
    A_j(S):=\sum_{i\in S}c_iw_{i,j}^{(1)}h_i^{\mathrm{in}},
    \qquad
    D_j(S):=\sum_{i\in S}c_iw_{i,j}^{(1)}.
    \label{eq:cache_additive_statistics}
\end{equation}
Because the summands are pointwise, their values on $S_0$ are unchanged after adding $B$, and
\begin{equation}
    A_j(S_1)=A_j(S_0)+A_j(B),
    \qquad
    D_j(S_1)=D_j(S_0)+D_j(B).
    \label{eq:cache_statistics_additivity}
\end{equation}
so the coarse values can be cached and only the contributions from $B$ need to be computed. We train a lightweight head $H_{\phi}$ with the frozen backbone, attach it to the first-block feature $h_i^{(1)}$, and define
\begin{equation}
    \widehat{u}^{\mathrm{head}}_i=H_{\phi}\!\left(h_i^{(1)}\right),
    \qquad
    s_{\mathrm{head}}(x_i)
    :=\Phi\!\left(
        \left\|\nabla\widehat{u}^{\mathrm{head}}(x_i)\right\|_2,
        \left\|\Delta\widehat{u}^{\mathrm{head}}(x_i)\right\|_2
    \right),
    \label{eq:learned_refinement_score}
\end{equation}
and use this score in \eqref{eq:unified_refinement_measure}. The head can alternatively learn $s_{\mathrm{head}}$ directly. Token attention, deslicing, and later blocks still run on $S_1$, because $B$ changes the global token states. Appendix~\ref{app:density_policy_design} gives the task-specific features and density constructions for all three policies.
